% answers_short.tex —— 2020 年·天津卷 · 简短答案速查
% 由 tools/make_short_answers.py 自 answers.tex 的 \gksolhead/\gkanswer 提取，
% 并与 2020/2020_数学真题_天津_答案要点.md 逐题比对后生成，勿手改（改源 MD 后重跑）。

\begin{gknotice}
\textbf{使用说明}：本节只给最终答案便于快速核对，完整推导见解析版。答案为编写者独立推导，\textbf{非官方标准答案}。
\end{gknotice}

\gkpart{选择题}
\begin{center}\begin{tabular}{cccccccccc}
\toprule
题号 & 1 & 2 & 3 & 4 & 5 & 6 & 7 & 8 & 9 \\
\midrule
答案 & C & A & A & B & C & D & D & B & D \\
\bottomrule
\end{tabular}\end{center}

\gkpart{填空题}
\begin{itemize}[leftmargin=2.6em]
  \item[\textbf{10.}] $3-2\mathrm i$。$\dfrac{8-\mathrm i}{2+\mathrm i}=\dfrac{(8-\mathrm i)(2-\mathrm i)}5=\dfrac{15-10\mathrm i}5=3-2\mathrm i$。
  \item[\textbf{11.}] $10$。通项为 $\mathrm C_5^kx^{5-3k}2^k$，令 $5-3k=2$ 得 $k=1$，系数为 $\mathrm C_5^1\cdot2=10$。
  \item[\textbf{12.}] $5$。圆心到直线的距离为 $\dfrac8{\sqrt{1+3}}=4$，半弦长为 $3$，故 $r=\sqrt{4^2+3^2}=5$。
  \item[\textbf{13.}] $\dfrac16$；$\dfrac23$。相互独立，都落入为 $\dfrac12\cdot\dfrac13=\dfrac16$，至少一个为 $1-\dfrac12\cdot\dfrac23=\dfrac23$。
  \item[\textbf{14.}] $4$。由 $ab=1$ 得 $\dfrac1{2a}+\dfrac1{2b}=\dfrac{a+b}2$，令 $t=a+b\geqslant2$，则原式为 $\dfrac t2+\dfrac8t\geqslant2\sqrt4=4$，$t=4$ 时取等。
  \item[\textbf{15.}] $\dfrac16$；$\dfrac{13}2$。以 $B$ 为原点、$BC$ 为 $x$ 轴建系得 $A\left(\dfrac32,\dfrac{3\sqrt3}2\right)$，由 $\overrightarrow{AD}\cdot\overrightarrow{AB}=-9\lambda=-\dfrac32$ 得 $\lambda=\dfrac16$；再设 $M(m,0)$，则 $\overrightarrow{DM}\cdot\overrightarrow{DN}=(m-2)^2+\dfrac{13}2\geqslant\dfrac{13}2$。
\end{itemize}

\gkpart{解答题}
\begin{itemize}[leftmargin=2.6em]
  \item[\textbf{16.}] (1) $C=\dfrac\pi4$。由 $\cos C=\dfrac{a^2+b^2-c^2}{2ab}=\dfrac{\sqrt2}2$ 得 $C=\dfrac\pi4$；(2) $\sin A=\dfrac{2\sqrt{13}}{13}$。由正弦定理 $\sin A=\dfrac{a\sin C}{c}=\dfrac{2\sqrt2\cdot\frac{\sqrt2}2}{\sqrt{13}}=\dfrac2{\sqrt{13}}$；(3) $\sin\left(2A+\dfrac\pi4\right)=\dfrac{17\sqrt2}{26}$。由 $\cos A=\dfrac3{\sqrt{13}}$ 得 $\sin2A=\dfrac{12}{13}$、$\cos2A=\dfrac5{13}$，故原式为 $\dfrac{17}{13\sqrt2}=\dfrac{17\sqrt2}{26}$
  \item[\textbf{17.}] (1) 证明。以 $C$ 为原点、$CA$、$CB$、$CC_1$ 为三轴建系，得 $\overrightarrow{C_1M}=(1,1,0)$、$\overrightarrow{B_1D}=(2,-2,-2)$，数量积为 $0$；(2) $\dfrac{\sqrt{30}}6$。平面 $BB_1E$ 即 $x=0$，平面 $DB_1E$ 的法向量为 $(1,-1,2)$，两法向量夹角余弦为 $\dfrac1{\sqrt6}$，故二面角正弦为 $\sqrt{1-\dfrac16}=\dfrac{\sqrt{30}}6$；(3) $\dfrac{\sqrt3}3$。$\overrightarrow{AB}=(-2,2,0)$ 与平面 $DB_1E$ 的法向量 $(1,-1,2)$ 夹角的余弦绝对值为 $\dfrac{|{-2-2}|}{2\sqrt2\cdot\sqrt6}=\dfrac{1}{\sqrt3}$，直线与平面所成角的正弦值即为此值
  \item[\textbf{18.}] (1) $\dfrac{x^2}{18}+\dfrac{y^2}9=1$。顶点 $A(0,-3)$ 给出 $b=3$，$|OF|=|OA|=3$ 给出 $c=3$，故 $a^2=b^2+c^2=18$；(2) $y=x-3$ 或 $y=\dfrac12x-3$。设 $AB:y=kx-3$，$B(u,ku-3)$ 在椭圆上给出 $u=\dfrac{12k}{1+2k^2}$，$CP\perp AB$ 给出 $u=\dfrac{2+6k}{1+k^2}$，联立得 $(2k-1)(k-1)=0$
  \item[\textbf{19.}] (1) $a_n=n$，$b_n=2^{n-1}$。由 $a_5=5(a_4-a_3)$ 得公差 $d=1$；由 $b_5=4(b_4-b_3)$ 得 $q^2=4(q-1)$，即 $(q-2)^2=0$，$q=2$；(2) 证明。$S_n=\dfrac{n(n+1)}2$，不等式两边约去 $\dfrac{(n+1)(n+2)}4$ 后化为 $n(n+3)<(n+1)(n+2)$，即 $0<2$；(3) $\dfrac{4^n}{2n+1}-\dfrac{6n+5}{9\cdot4^n}-\dfrac49$。奇数项满足 $\dfrac{(6j-5)4^{j-1}}{(2j-1)(2j+1)}=\dfrac{4^j}{2j+1}-\dfrac{4^{j-1}}{2j-1}$ 可裂项，偶数项为等比与等差之积，分组求和
  \item[\textbf{20.}] (1) ① $y=9x-8$；② 单调递减区间 $(0,1)$，单调递增区间 $(1,+\infty)$，极小值 $g(1)=1$，无极大值。$f(1)=1$、$f'(1)=9$ 给出切线；$g'(x)=\dfrac{3(x-1)^3(x+1)}{x^2}$ 定号即得单调性与极值；(2) 证明。令 $F(x)=f(x)-f(x_2)-\dfrac{x-x_2}2\left[f'(x)+f'(x_2)\right]$，则 $F(x_2)=0$ 且 $F'(x)=-\dfrac{(x-x_2)^2}2\left(3+\dfrac{k}{x^2x_2}\right)$；由 $x\geqslant x_2\geqslant1$、$k\geqslant-3$ 得括号内非负，故 $F$ 严格递减，$F(x_1)<0$ 即所证
\end{itemize}
